% Sections 15.1 to 15.6, translated from sv/chapters/15/theory.tex.

\section{Introduction}
\label{c15:sec:inledning}
Complex numbers are needed when we want to solve equations such as $x^2 = -1$. We introduce the
\textbf{imaginary unit}~$j$:
\[
  j^2 = -1 \quad \Rightarrow \quad j = \sqrt{-1}
\]
In electrical engineering, $j$ is used instead of $i$, which is reserved for current.

\section{Rectangular form}
\label{c15:sec:rektangular}
A complex number $z$ is written in \textbf{rectangular form} as
\[
  z = x + jy,
\]
where
\begin{itemize}
  \item $x = \operatorname{Re}(z)$ is the real part,
  \item $y = \operatorname{Im}(z)$ is the imaginary part.
\end{itemize}
Complex numbers are drawn in the \textbf{complex plane}, with the Re axis horizontal and the Im axis
vertical.

\section{Magnitude and phase angle}
\label{c15:sec:fasvinkel}
The \textbf{magnitude}, or absolute value, is the distance from the origin:
\[
  \abs{z} = \sqrt{x^2 + y^2}
\]
The \textbf{phase angle} is the angle to the positive Re axis:
\[
  \delta = \arctan\left(\frac{y}{x}\right) + \text{quadrant correction}
\]
The correction depends on which quadrant the number lies in:

\begin{narrowtable}{@{}l@{\hspace{2.5em}}c@{\hspace{2.5em}}c@{\hspace{2.5em}}l@{}}
  \thead{Quadrant} & \thead{$x$} & \thead{$y$} & \thead{Correction} \\
  \midrule
  I & $+$ & $+$ & none \\
  II & $-$ & $+$ & $+180^{\circ}$ \enspace ($+\pi$) \\
  III & $-$ & $-$ & $+180^{\circ}$ \enspace ($+\pi$) \\
  IV & $+$ & $-$ & none (negative $\delta$) \\
\end{narrowtable}

\section{Polar form}
\label{c15:sec:polar}
In \textbf{polar form}, a complex number is written with its magnitude and its phase angle:
\[
  z = \abs{z}\angle\delta
\]
The conversion from \textbf{rectangular to polar form} is given by
\[
  \abs{z} = \sqrt{x^2 + y^2}, \qquad \delta = \arctan(y/x) + \text{correction},
\]
and the conversion from \textbf{polar to rectangular form} by
\[
  x = \abs{z}\cos\delta, \qquad y = \abs{z}\sin\delta.
\]

\section{Operations in rectangular form}
\label{c15:sec:rakneoperationer}
\textbf{Addition and subtraction:} the real and imaginary parts are added or subtracted
separately.
\[
  (x_1 + jy_1) \pm (x_2 + jy_2) = (x_1 \pm x_2) + j(y_1 \pm y_2)
\]
\textbf{Multiplication:} multiply out and use $j^2 = -1$.
\[
  (x_1 + jy_1)(x_2 + jy_2) = (x_1x_2 - y_1y_2) + j(x_1y_2 + x_2y_1)
\]
\textbf{Conjugate:} the conjugate of $z = x + jy$ is $\bar{z} = x - jy$. It is useful in
division, since it makes the denominator real.

\section{Worked examples}
\label{c15:sec:typexempel}

\begin{exempel}[Complex numbers in the complex plane]
Mark $z_1 = 3 + j4$ and $z_2 = -2 - j3$ in the complex plane and calculate their magnitudes and
phase angles.

The number $z_1 = 3 + j4$ lies in quadrant I:
\[
  \abs{z_1} = \sqrt{9 + 16} = 5, \qquad
  \delta_1 = \arctan\left(\frac{4}{3}\right) \approx 53.1^{\circ}
\]
In polar form, $z_1 = 5\angle 53.1^{\circ}$.

The number $z_2 = -2 - j3$ lies in quadrant III, so $180^{\circ}$ is added:
\begin{gather*}
  \abs{z_2} = \sqrt{4 + 9} = \sqrt{13} \approx 3.61, \\
  \delta_2 = \arctan\left(\frac{-3}{-2}\right) + 180^{\circ} \approx 56.3^{\circ} + 180^{\circ}
    = 236.3^{\circ}
\end{gather*}
\end{exempel}

\begin{exempel}[Converting from rectangular to polar form]
The voltage is $U = 3 + j4\,\text{V}$. Find its magnitude and phase angle.
\[
  \abs{U} = \sqrt{3^2 + 4^2} = 5\,\text{V}, \qquad
  \delta = \arctan\left(\frac{4}{3}\right) \approx 0.927\,\text{rad} \approx 53.1^{\circ}
\]
In polar form, $U = 5\angle 0.927\,\text{rad}$.
\end{exempel}

\needspace{9\baselineskip}
\begin{exempel}[Converting from polar to rectangular form]
The current is $I = 10\angle\frac{\pi}{4}\,\text{mA}$. Find its rectangular form.
\begin{gather*}
  x = 10\cos\left(\frac{\pi}{4}\right) = \frac{10}{\sqrt{2}} \approx 7.07\,\text{mA} \\
  y = 10\sin\left(\frac{\pi}{4}\right) = \frac{10}{\sqrt{2}} \approx 7.07\,\text{mA} \\
  I = (7.07 + j7.07)\,\text{mA}
\end{gather*}
\end{exempel}
